
%=====================================================================
\section{Mathematical significance and consequences}\label{sec:signif}
%=====================================================================
\subsection{What has been established}
Table~\ref{tab:results} collects the results with their scope and status. The central statement is Theorem~\ref{thm:B}: the autonomous rule \textsf{MAXPP}, which consults nothing but the integer it is acting on, generates from \emph{every} starting integer $x_0\ge2$ an itinerary that is not eventually periodic and whose eventual factor complexity is exactly $k+1$. Everything else in the paper either explains why this holds (Theorem~\ref{thm:A}, Sections~\ref{sec:fire}--\ref{sec:rotation}), sharpens it (Theorem~\ref{thm:C}), or delimits it (Theorems~\ref{thm:D}, \ref{thm:E}, \ref{thm:ordE}).
\begin{table}[H]\centering\small
\begin{tabular}{@{}p{0.15\textwidth}p{0.38\textwidth}p{0.2\textwidth}p{0.15\textwidth}@{}}\toprule
Result & Statement & Scope & Status\\\midrule
Thm.~\ref{thm:A} & primes $\ge5$ are fired finitely often; tail alphabet $\{2,3\}$ & every prime-fire selector, every $x_0$ & proved\\
Thms.~\ref{thm:spf}, \ref{thm:lpf} & \textsf{SPF}: $(3,2,2,2,2)^\infty$; \textsf{LPF}: $(2,3,3)^\infty$ & every $x_0$ & proved\\
\textsf{MINPP} & eventual period $3$ or $5$ & $2\le x_0\le400$ & computed\\
Thm.~\ref{thm:B} & \textsf{MAXPP}: not eventually periodic, $p(k)=k+1$, $h=0$ & every $x_0\ge2$ & proved\\
Thm.~\ref{thm:C} & frequency $\ln48/\ln864$, discrepancy $<1$, growth rate $2.2747033$ bits/step & every $x_0\ge2$ & proved\\
Thm.~\ref{thm:D} & finite prime support $\iff$ finitely many branch values & every selector & proved (Størmer imported)\\
Thm.~\ref{thm:E} & $|W_k|\le|Q|M^dc_dn_k^d$; entropy $0$ for all finite codings of branch/control & finite-guard class & proved\\
Thms.~\ref{thm:ordS}, \ref{thm:ordE} & counter selector; $h(u)\ge\frac2{\pi^2}\ln2$ & the counter $S$ & proved\\
\bottomrule\end{tabular}
\caption{Results, scope and status.}\label{tab:results}\end{table}

\subsection{One irrational number}
The transition grammar turns each branch into an integer translation of the valuation vector, and the selector \textsf{MAXPP} turns the choice between $2$ and $3$ into the sign of a single real number $f=\alpha\ln2-\beta\ln3$. The two translations change $f$ by $-\ln18$ and $+\ln48$, whose sum $\ln864$ is the length of the circle on which the process is a rotation. The rotation number $\theta=\ln48/\ln864$ is irrational because $48^s=864^r$ is impossible, and that single fact gives non-periodicity, the Sturmian complexity $k+1$, the equidistribution of the phases, the exact frequency of each prime, the bounded discrepancy, and the Beatty form of the valuations. The only arithmetic input is unique factorisation. The result is therefore a theorem about integer arithmetic, not an experimental regularity, and it holds for the entire orbit space $x_0\ge2$ at once.

Three features distinguish the situation from a rotation chosen by hand. First, the rotation is \emph{generated}: no external clock, input or memory enters; the angle is fixed by the grammar (the exponents of the two prime branches), and the starting integer fixes only the phase. Second, the underlying integer state grows without bound, $\log_2x_n\ge2n$, so the Sturmian itinerary is produced by an unbounded orbit and not by a bounded compact system. Third, the quantitative structure is exact: Theorem~\ref{thm:C} gives the frequency, the discrepancy, the Beatty form and the growth rate in closed form.

\subsection{A dichotomy between finite and unbounded branching}
Theorems~\ref{thm:D} and~\ref{thm:E} describe the space of selectors from two sides. If only finitely many branch values occur, the total prime support is finite and the dynamics admits a finite-dimensional prime-valuation representation as integer translations. Within the class of selectors whose decisions are determined by finite affine valuation guards, thresholds, congruences and finite control, Theorem~\ref{thm:E} further gives polynomial branch/control factor complexity and zero entropy; \textsf{MAXPP} lies at the bottom of the unbounded part of this regime, with complexity $k+1$. If infinitely many branch values occur, the prime support renews, and the dynamics leaves the finite-dimensional regime altogether. The counter of Section~\ref{sec:order} is the simplest member of this second class; it generates nothing more complicated than $j\mapsto j+1$, yet it carries an observable whose entropy is positive. This is why the paper distinguishes the complexity \emph{generated} by selection from the complexity \emph{revealed} by an observable: the first is a property of the dynamics, the second of the arithmetic of the values that the dynamics enumerates.

\subsection{Consequences}
\begin{enumerate}[label=(\arabic*),itemsep=2pt]
\item \textbf{Autonomous aperiodicity holds for every initial integer.} For the magnitude rule \textsf{MAXPP} the non-periodic itinerary does not depend on a special starting point or on a designated orbit; every $x_0\ge2$ enters the same rotation after at most a transient (Figures~\ref{fig:plane}, \ref{fig:transients}), and only the phase $z_0$ differs.
\item \textbf{Prime-fire dynamics is two-dimensional in the limit.} By Theorem~\ref{thm:A} the infinite-dimensional valuation lattice collapses, for every prime-fire selector, to the plane $(v_2,v_3)$ with the two translations $(-1,+2)$ and $(+4,-1)$. The three elementary rules \textsf{SPF}, \textsf{LPF}, \textsf{MAXPP} then differ only in the rule that picks between two translations: a periodic rule, a different periodic rule, and a comparison of magnitudes that is a rotation.
\item \textbf{Zero entropy for finite-guard selectors.} Finite branch sets imply finite total prime support, but not by themselves zero entropy. Within the class of selectors whose decisions are determined by finite affine valuation guards, thresholds, congruences and finite control, Theorem~\ref{thm:E} gives polynomial growth of the number of branch/control words and zero entropy; positive entropy of a branch/control itinerary in this setting therefore requires infinitely many branch values or guards outside the class.
\item \textbf{Observables must be read with their source.} A finite-valued function of an unbounded branch value can have positive entropy while the selection dynamics is a counter. The criterion ``a fixed finite observable of the branch itinerary has positive entropy'' therefore does not by itself certify complex selection; it certifies complexity of the observable's arithmetic.
\item \textbf{Exact growth of the substrate.} The rate $2.2747033\ldots$ bits per step, the Beatty form $\alpha_n,\beta_n$, and the residual bound $\log_2\frac{32}{27}$ describe how the integer state itself grows along the Sturmian itinerary.
\end{enumerate}

\subsection{Outlook within the stated scope}
The statements above are about the forward transition $T_N$, the selectors defined in this paper, and the observables defined here. The proofs use only unique factorisation, the Chinese remainder theorem, Euler's product, and the classical results cited in the appendices.

%=====================================================================
\appendix
\section{Classical lemmas with proofs}\label{app:classical}
%=====================================================================
\begin{lemma}[Morse--Hedlund]\label{lem:mh}
Let $w$ be an infinite one-sided word over a finite alphabet. If $p_w(k)\le k$ for some $k\ge1$, then $w$ is eventually periodic. Equivalently, a word that is not eventually periodic has $p_w(k)\ge k+1$ for all $k\ge1$.
\end{lemma}
\begin{proof}
Every factor of length $j$ has at least one right extension in $w$ (the word is infinite), so $p_w(j+1)\ge p_w(j)$. If $p_w(1)=1$ the word is constant. Otherwise $p_w(1)\ge2$. If $p_w(j+1)>p_w(j)$ for all $j<k$, then $p_w(k)\ge p_w(1)+(k-1)\ge k+1$, contrary to the hypothesis. Hence $p_w(j)=p_w(j+1)$ for some $j<k$. The map sending a factor of length $j+1$ to its prefix of length $j$ is onto the set of factors of length $j$ and the two sets have equal size, so it is a bijection: each factor $u$ of length $j$ has exactly one right extension, i.e.\ $w_{i+j}=F(w_i\cdots w_{i+j-1})$ for all $i\ge0$ for a function $F$. The sequence of windows $W_i=w_i\cdots w_{i+j-1}$ is then the orbit of a map on the finite set of factors of length $j$, hence eventually periodic, and $w$ is eventually periodic. \end{proof}

\begin{proof}[Proof of Lemma~\ref{lem:regions}]
Let $R(n,d)$ be the maximum number of regions of $n$ distinct affine hyperplanes in $\R^d$. Then $R(0,d)=1$ and $R(n,0)=1$. For $n\ge1$ remove $H_n$. The $n-1$ remaining hyperplanes give at most $R(n-1,d)$ regions. Adding $H_n$ splits each region that $H_n$ meets into exactly two convex pieces, and the regions that $H_n$ meets correspond bijectively to the regions into which the traces $H_i\cap H_n$ ($i<n$) cut $H_n\cong\R^{d-1}$; there are at most $R(n-1,d-1)$ of them. Hence $R(n,d)\le R(n-1,d)+R(n-1,d-1)$, and induction with Pascal's rule gives $R(n,d)\le\sum_{i\le d}\binom ni$. Repeated hyperplanes only reduce the count. \end{proof}

\begin{proof}[Proof of Lemma~\ref{lem:faces}]
Let $\sigma$ be a realised sign vector and $I=\{i:\sigma_i=0\}$. The set $F_I=\{y:g_i(y)=0,\ i\in I\}$ is a nonempty affine subspace; choose canonically (e.g.\ lexicographically first) a subset $J\subset I$ whose linear parts form a basis of the span of the linear parts of $\{g_i\}_{i\in I}$. Since the system is consistent, $F_I=F_J$, and $\dim F_J=d-|J|$. A point $y$ with sign vector $\sigma$ lies in $F_J$, and each $g_i$ with $i\notin I$ is nonzero at $y$. The functions $g_i|_{F_J}$ that are not identically zero cut $F_J$ into regions; $y$ lies in one of them and $\sigma$ is determined by $J$ and this region (the functions identically zero on $F_J$ have sign $0$, the others a constant nonzero sign on the region). Hence the number of sign vectors is at most
\[
\sum_{j=0}^{d}\binom nj\sum_{i=0}^{d-j}\binom ni\ \le\ \sum_{j=0}^{d}\sum_{i=0}^{d-j}n^{j+i}\ \le\ \frac{(d+1)(d+2)}2\,n^d,
\]
using Lemma~\ref{lem:regions} on $F_J\cong\R^{d-j}$ (at most $n$ nonconstant traces), $\binom nm\le n^m$, and $n\ge1$. \end{proof}

\paragraph{Imported results.} Størmer's theorem (Theorem~\ref{thm:stormer}) and Euler's product $\prod_{l\ \mathrm{prime}}(1-l^{-2})=6/\pi^2$ are used without proof.

%=====================================================================
\section{Exact certificates and computational reproductions}\label{app:checks}
%=====================================================================
The script \texttt{code/verify\_prime\_dyn.py} implements two independent engines (exact prime-exponent dictionaries and big-integer arithmetic) and the checks below; it reports \texttt{ALL PASS} and writes \texttt{reports/verification\_report.json}. Integer identities are compared exactly; real quantities use $60$-digit arithmetic with an explicit margin. Computation is never used to prove an asymptotic claim; each check tests a proved statement or documents a table.
\begin{table}[ht]\centering\small
\begin{tabular}{@{}ll@{}}\toprule
Statement & Checks\\\midrule
Valuation law (Thm.~\ref{thm:valuation}) & C1, C2\\
Prime-fire reduction (Thm.~\ref{thm:A}), \textsf{SPF}, \textsf{LPF} & C3, C7a, C7b\\
Two-counter law, window, blocks (Prop.~\ref{prop:twocounter}, Lem.~\ref{lem:window}, Prop.~\ref{prop:blocks}) & C4a, C4b\\
Complexity $k+1$, non-periodicity (Thms.~\ref{thm:cx}, \ref{thm:nonper}) & C4c, C4d\\
Mechanical word, discrepancy, $\theta$ (Cor.~\ref{cor:mech}, Thm.~\ref{thm:C}) & C5a, C5b, C5c\\
Beatty valuations (Cor.~\ref{cor:beatty}) & C5d\\
Step ratios and growth rate (Thm.~\ref{thm:C}(iii)) & C6a, C6b\\
Størmer tables, renewal lemma, measurements (§\ref{sec:composite}) & C8a, C8b, C9\\
Face determinism and polynomial bound (Thm.~\ref{thm:E}) & C10a, C10b\\
Counter, $u$, CRT construction, density, block counts (§\ref{sec:order}) & C11a--C11g\\
Inherited archives (hash check only) & C12\\
\bottomrule\end{tabular}
\caption{Where each proved statement is exercised by the verifier.}\label{tab:map}\end{table}
\begin{center}\small{\small\begin{longtable}{@{}lr@{}}\toprule
check & outcome\\\midrule\endhead
\texttt{C1\_valuation\_law} & pass \\
\texttt{C2\_exponent\_vs\_bigint\_implementations} & pass \\
\texttt{C3\_prime\_fire\_reduction} & pass \\
\texttt{C4a\_two\_counter\_law\_and\_regime\_entry} & pass \\
\texttt{C4b\_tail\_has\_no\_33\_and\_no\_222} & pass \\
\texttt{C4c\_factor\_complexity\_k\_plus\_1} & pass \\
\texttt{C4d\_no\_period\_up\_to\_1500} & pass \\
\texttt{C5a\_mechanical\_word\_formula} & pass \\
\texttt{C5b\_discrepancy\_below\_one} & pass \\
\texttt{C5c\_theta\_value} & pass \\
\texttt{C5d\_closed\_form\_valuations} & pass \\
\texttt{C6a\_step\_ratios\_9/2\_and\_16/3} & pass \\
\texttt{C6b\_growth\_rate} & pass \\
\texttt{C7a\_SPF\_period5} & pass \\
\texttt{C7b\_LPF\_period3} & pass \\
\texttt{C8a\_stormer\_tables} & pass \\
\texttt{C8b\_renewal\_lemma} & pass \\
\texttt{C9\_renewal\_measurements} & pass \\
\texttt{C10a\_same\_face\_same\_residue\_same\_word} & pass \\
\texttt{C10b\_polynomial\_count\_bound} & pass \\
\texttt{C11a\_decoder\_selector} & pass \\
\texttt{C11b\_u\_zero\_iff\_equal\_exponents} & pass \\
\texttt{C11g\_N\_equals\_x\_capped\_valuations\_periodic} & pass \\
\texttt{C11c\_self\_referential\_N\_equals\_x} & pass \\
\texttt{C11d\_crt\_construction} & pass \\
\texttt{C11e\_mean\_density\_bound} & pass \\
\texttt{C11f\_block\_counts\_u} & pass \\
\texttt{C12\_inherited\_archive\_hashes} & pass \\
\bottomrule\end{longtable}}
\end{center}
All checks pass in the delivered release (about $75$ seconds on the build machine).

%=====================================================================
\section{Source provenance and inherited material}\label{app:prov}
%=====================================================================
The selector \textsf{MAXPP}, the prime-fire reduction, the exact frequencies, the finite-guard counting argument and the counter selector were first formulated and tested in three earlier releases of the UNNS programme, which are preserved unchanged and enter this release as inputs:
\begin{center}\scriptsize
\begin{tabular}{@{}ll@{}}\toprule
archive & SHA-256\\\midrule
\texttt{UNNS\_U4\_GEN\_v001.zip} & \texttt{042eb5b09b4d231283c9c7d9d40e29a8789625f4b073e724000d49f6d49e18df}\\
\texttt{UNNS\_U4\_GEN\_C\_v001.zip} & \texttt{7a361885e66c2448c232dc04b3c6c52c92a3cae7ec4452ed5dfd63fdade59118}\\
\texttt{UNNS\_U4\_GEN\_C\_ORDER\_v001.zip} & \texttt{a50227758ce1d5f819ae6a1775de2e9dd1563cc7a283b86244b07ed27dceb975}\\
\bottomrule\end{tabular}\end{center}
The verifiers of these archives are \textsc{NOT\_RERUN} here; check C12 verifies only that the archives are present and bit-identical to the hashes above. All mathematics stated in this paper is proved in the text or in the appendices and is checked by the independent verifier of this release; the earlier releases are not relied upon for any step of any proof. In particular the polynomial bound of Theorem~\ref{thm:E} is the bound proved in Section~\ref{sec:complexity}, and no quantitative bound from an earlier release is used.

\paragraph{Where the sources are held.} \textsf{U4-GEN} is part of the commit-pinned collection~\cite{coll}; its release directory and proof document are given in~\cite{u4gen}. \textsf{U4-GEN-C} and \textsf{U4-GEN-C-ORDER} are \emph{not} included in that collection. They are identified here only by release name and archive SHA-256 (table above), and the archives are preserved unchanged in \texttt{inherited/}; no statement is claimed to be found at a GitHub location for them. \textsf{U4-GEN-C-ORDER} also records the withdrawal of the quantitative bound GEN-C3a of \textsf{U4-GEN-C} (ORD-0, E1); that bound is not used in this paper.

\paragraph{Source of each theorem family.} ``Inherited'' means that the statement was formulated in the named release; every such statement is proved in the text, and no proof relies on the release. ``Derived here'' means that no release is the source. ``Recorded verification'' refers to the checks of Appendix~\ref{app:checks} run by the verifier of this release.
\begin{center}\scriptsize
\begin{tabular}{@{}p{0.30\textwidth}p{0.36\textwidth}p{0.26\textwidth}@{}}\toprule
Result & Source release and item & Kind\\\midrule
Valuation law, legality and growth (Thm.~\ref{thm:valuation}, Prop.~\ref{prop:growth}) & none (grammar of the transition) & derived here\\
Prime-fire reduction (Thm.~\ref{thm:A}) & U4-GEN, GEN-2 & inherited, re-proved\\
\textsf{SPF} (Thm.~\ref{thm:spf}) & U4-GEN, GEN-3 & inherited, re-proved\\
\textsf{LPF} (Thm.~\ref{thm:lpf}) & U4-GEN, registered observation only & derived here\\
\textsf{MINPP} (observed) & U4-GEN, registered observation & recorded verification\\
\textsf{MAXPP}: non-periodicity, $p(k)=k+1$, $h=0$ (Thm.~\ref{thm:B}) & U4-GEN, GEN-4 (Sturmian theorem imported) & inherited, re-proved\\
Frequencies, discrepancy, growth (Thm.~\ref{thm:C}) & U4-GEN-C, GEN-5 & inherited, re-proved\\
Support renewal (Thm.~\ref{thm:D}) & U4-GEN-C, Prop.~S (Størmer imported) & inherited, re-proved\\
Finite-guard complexity (Thm.~\ref{thm:E}) & U4-GEN-C, GEN-C3 (GEN-C3a withdrawn, not used) & inherited, re-proved\\
Counter selector and exponent profile (Thms.~\ref{thm:ordS}, \ref{thm:ordE}, Prop.~\ref{prop:ord2}) & U4-GEN-C-ORDER, ORD-2, ORD-3, ORD-E, ORD-C & inherited, re-proved\\
Tables, figures, checks C1--C12 & this release (\texttt{code/verify\_prime\_dyn.py}) & recorded verification\\
\bottomrule\end{tabular}\end{center}

\paragraph{Classical sources.} Fekete's lemma; Størmer's theorem; Euler's product for $\zeta(2)$; the Morse--Hedlund characterisation of Sturmian-like complexity; the region and face counts for hyperplane arrangements (Buck, Zaslavsky). They are cited in the bibliography and proved or credited where used.

%=====================================================================
\section{Package contents and reproduction}\label{app:pkg}
%=====================================================================
The release \texttt{UNNS\_PRIME\_DYN\_v003} contains \texttt{code/} (engine \texttt{pdyn.py}, verifier, figure and table generators), \texttt{data/} (computed data files), \texttt{figures/} and \texttt{tables/} (generated), \texttt{manuscript/} (the \LaTeX{} sources), \texttt{reports/} (the verification report and its log), \texttt{inherited/} (the three archives and their hashes), \texttt{README.md}, \texttt{SHA256SUMS.txt} and \texttt{run\_all.sh}. Running \texttt{run\_all.sh} executes the verifier, regenerates the figures and tables, and rebuilds the PDF. The compiled manuscript is Version~1.2.

\begin{thebibliography}{9}\small
\emergencystretch=4em
\bibitem{coll} UNNS Research Collective, \emph{UNNS Arithmetic Dynamics} (curated collection of the releases U1--U4), commit \texttt{22d8c2e23f49fb98a5b04c6b7efe9ba0a3703808}. \url{https://github.com/ukbbi/UNNS/tree/22d8c2e23f49fb98a5b04c6b7efe9ba0a3703808/projects/UNNS_ARITHMETIC_DYNAMICS}.
\bibitem{u4gen} UNNS Research Collective, \emph{Intrinsic Branching and Autonomous Generative Dynamics}, UNNS-U4-GEN, version v001 (no manuscript; proofs in \texttt{docs/PROOFS\_GEN.md}). Release directory: \url{https://github.com/ukbbi/UNNS/tree/22d8c2e23f49fb98a5b04c6b7efe9ba0a3703808/projects/UNNS_ARITHMETIC_DYNAMICS/04_symbolic_dynamics/UNNS_U4_GEN_v001}; proofs: \url{https://github.com/ukbbi/UNNS/blob/22d8c2e23f49fb98a5b04c6b7efe9ba0a3703808/projects/UNNS_ARITHMETIC_DYNAMICS/04_symbolic_dynamics/UNNS_U4_GEN_v001/docs/PROOFS_GEN.md}. Archive SHA-256: \path{042eb5b09b4d231283c9c7d9d40e29a8789625f4b073e724000d49f6d49e18df}.
\bibitem{u4genc} UNNS Research Collective, \emph{Composite-Branch Dynamics and Autonomous Entropy}, UNNS-U4-GEN-C, version v001 (archive \texttt{UNNS\_U4\_GEN\_C\_v001.zip}). Not included in the commit-pinned collection; identified by archive SHA-256 \path{7a361885e66c2448c232dc04b3c6c52c92a3cae7ec4452ed5dfd63fdade59118}; the archive is preserved in \texttt{inherited/} of this release.
\bibitem{u4order} UNNS Research Collective, \emph{Valuation-Order Observables}, UNNS-U4-GEN-C-ORDER, version v001 (archive \texttt{UNNS\_U4\_GEN\_C\_ORDER\_v001.zip}). Not included in the commit-pinned collection; identified by archive SHA-256 \path{a50227758ce1d5f819ae6a1775de2e9dd1563cc7a283b86244b07ed27dceb975}; the archive is preserved in \texttt{inherited/} of this release.
\bibitem{fekete} M.\ Fekete, \emph{\"Uber die Verteilung der Wurzeln bei gewissen algebraischen Gleichungen mit ganzzahligen Koeffizienten}, Math.\ Z.\ 17 (1923), 228--249.
\bibitem{stormer} C.\ Størmer, \emph{Quelques théorèmes sur l'équation $\mathrm{Arctan}\frac1x+\mathrm{Arctan}\frac1y=\frac k4\pi$}, Nouv.\ Ann.\ Math.\ (1897), and related papers on consecutive smooth numbers.
\bibitem{mh} M.\ Morse and G.\ A.\ Hedlund, \emph{Symbolic dynamics II. Sturmian trajectories}, Amer.\ J.\ Math.\ 62 (1940), 1--42.
\bibitem{lothaire} M.\ Lothaire, \emph{Algebraic Combinatorics on Words}, Cambridge University Press, 2002 (Chapter 2: Sturmian words).
\bibitem{buck} R.\ C.\ Buck, \emph{Partition of space}, Amer.\ Math.\ Monthly 50 (1943), 541--544.
\bibitem{zaslavsky} T.\ Zaslavsky, \emph{Facing up to arrangements: face-count formulas for partitions of space by hyperplanes}, Mem.\ Amer.\ Math.\ Soc.\ 154 (1975).
\bibitem{sarnak} P.\ Sarnak, \emph{Three lectures on the Möbius function, randomness and dynamics}, IAS lecture notes, 2010 (squarefree integers as a dynamical system).
\end{thebibliography}
